% 当前读者版插图与地图来源。
\creditheading{封面}
封面采用英国画家 J.\,M.\,W. Turner 1815 年油画《狄多建造迦太基，或迦太基帝国的兴起》（\textit{Dido building Carthage, or The Rise of the Carthaginian Empire}），现藏伦敦国家美术馆，馆藏编号 NG498。所用高清图像来自 Wikimedia Commons，属公有领域。

\creditheading{卷首版画}
十二幅卷首画均为 Wenceslaus Hollar 据 Francis Cleyn 设计刻制的《埃涅阿斯纪》蚀刻版画，约作于 1654 年；扫描图像来自多伦多大学 Wenceslaus Hollar Digital Collection，并经 Wikimedia Commons 提供，均属公有领域。图版经裁切，以灰阶呈现。

\creditline{第一卷}{《狄多与埃涅阿斯》}
\creditline{第二卷}{《埃涅阿斯离开特洛伊》}
\creditline{第三卷}{《埃涅阿斯与安德罗玛刻相会》}
\creditline{第四卷}{《墨丘利催促埃涅阿斯》}
\creditline{第五卷}{《埃涅阿斯与阿刻斯忒斯》}
\creditline{第六卷}{《埃涅阿斯与卡戎》}
\creditline{第七卷}{《阿玛塔与拉维尼娅在酒神女信徒之中》}
\creditline{第八卷}{《维纳斯将武器交给埃涅阿斯》}
\creditline{第九卷}{《尼苏斯杀死沃尔肯斯》}
\creditline{第十卷}{《埃涅阿斯与墨泽提乌斯、劳苏斯之战》}
\creditline{第十一卷}{《卡米拉杀死奥努斯之子》}
\creditline{第十二卷}{《埃涅阿斯与图尔努斯的最后决斗》}

\creditheading{地图}
\creditline{地图 1}{斯特拉波所述的世界。}
\creditline{地图 2}{从特洛伊到拉提乌姆。}
\creditline{地图 3}{拉提乌姆与台伯河。}
\par\addvspace{3mm}\noindent
航程图与拉提乌姆、台伯河图据诗中叙事绘制，路线与战事范围为示意。古代岸线数据来自 Ancient World Mapping Center（AWMC，ODbL），现代岸线来自 Natural Earth（公有领域）。斯特拉波世界图采用 Bunbury 1883 年重构图，图版出自《古代地理史》（\textit{A History of Ancient Geography}）第二卷。
